% Reviewed source SHA256 (UTF-8/LF): 8c6e572e0a95b4ff198de3ed9403c007b90f51047f19ccbce3ebdebc3684dde8
\documentclass[11pt]{article}
\usepackage[T1]{fontenc}
\usepackage{lmodern}
\usepackage[a4paper,margin=27mm]{geometry}
\usepackage{amsmath,amssymb,amsthm,mathtools}
\usepackage{microtype,booktabs,array}
\usepackage[hidelinks]{hyperref}
\usepackage{enumitem}
\setlist{itemsep=3pt,topsep=5pt}
\newtheorem{theorem}{Theorem}[section]
\newtheorem{lemma}[theorem]{Lemma}
\newtheorem{corollary}[theorem]{Corollary}
\newtheorem{proposition}[theorem]{Proposition}
\theoremstyle{definition}
\newtheorem{definition}[theorem]{Definition}
\newtheorem{remark}[theorem]{Remark}
\newcommand{\R}{\mathbb R}
\newcommand{\E}{\mathbb E}
\newcommand{\Prob}{\mathbb P}
\newcommand{\tr}{\operatorname{tr}}
\newcommand{\rank}{\operatorname{rank}}
\newcommand{\diag}{\operatorname{diag}}
\newcommand{\range}{\operatorname{range}}
\newcommand{\OPT}{\operatorname{OPT}}
\newcommand{\norm}[1]{\left\lVert#1\right\rVert}
\newcommand{\ip}[2]{\left\langle#1,#2\right\rangle}
\newcommand{\psd}{\succeq}
\newcommand{\pd}{\succ}
\newcommand{\eps}{\varepsilon}
\DeclareMathOperator*{\argmin}{arg\,min}
\allowdisplaybreaks[1]
\hypersetup{pdfcreator={LaTeX},pdfauthor={Matthew J. Colbrook}}

\title{Sharp factor-two excess transfer for commuting positive semidefinite matrices}
\author{Matthew J. Colbrook\\\small Department of Applied Mathematics and Theoretical Physics\\\small University of Cambridge, Cambridge, United Kingdom\\\small\href{mailto:m.colbrook@damtp.cam.ac.uk}{m.colbrook@damtp.cam.ac.uk}}
\date{11 September 2026}
\usepackage{xurl}
\setlength{\emergencystretch}{3em}
\begin{document}
\maketitle

\noindent\textbf{Verified partial result for RA-10.} Theorem 1.1 and Section 3 give sharp factor-two transfer for a simultaneous eigenbasis and specified truncation projector. The general noncommuting nuclear question remains open.
Independent agent verification is not external human peer review or formal certification. Original draft remarks about review or source access reflect the input date; the dated \href{https://github.com/MColbrook/OpenProblemsInNLA/blob/codex/colbrook-transfer-resolutions/references/colbrook-transfer-2026-09-11/verification/reviews/RA-10-review.md}{independent review} records the subsequent checks.
\par\medskip
% BEGIN REVIEWED BODY

\begin{abstract}
For every finite Schatten norm, we prove a sharp factor-two bound on relative approximation excess when the positive semidefinite input, its approximation, and the selected truncation projector have a common orthonormal eigenbasis. This is a partial result for the nuclear norm: it does not resolve the corresponding question for arbitrary noncommuting positive semidefinite matrices.
\end{abstract}

\section{Statement}
Let $f:[0,\infty)\to[0,\infty)$ be continuous and nondecreasing, with $f(x)/x$ nonincreasing on $(0,\infty)$. No assumption that $f(0)=0$ is made.
Let $A,B\succeq0$ have a common orthonormal eigenbasis $q_1,\ldots,q_n$, ordered so that the eigenvalues of $A$ satisfy $a_1\ge\cdots\ge a_n\ge0$. Fix $1\le k<n$ and a set $S\subseteq\{1,\ldots,n\}$ of exactly $k$ indices, and suppose $b_i=0$ for $i\notin S$. Define
\[
 A=\sum_{i=1}^n a_iq_iq_i^T,\qquad
 B=\sum_{i\in S}b_iq_iq_i^T,\qquad
 P=\sum_{i\in S}q_iq_i^T,\qquad
 C=\sum_{i\in S}f(b_i)q_iq_i^T.
\]
In particular, the selected projector $P$ also shares this eigenbasis. The definition of $C$ includes $f(0)$ at the padded zero eigenvalues indexed by $S$; it does not mean the full matrix $f(B)$ when $f(0)>0$.

For $1\le p<\infty$, let $\norm{M}_{S_p}=(\sum_i\sigma_i(M)^p)^{1/p}$ and put
\[
 T_A=\sum_{i>k}a_i^p,\qquad
 T_f=\sum_{i>k}f(a_i)^p.
\]
These are the optimal rank-$k$ Schatten $p$-errors raised to the power $p$.

\begin{theorem}\label{thm:commuting}
Suppose $\tau=a_{k+1}>0$ and set $c=f(\tau)/\tau$. Then
\[
 \norm{f(A)-C}_{S_p}^p-T_f
 \le 2c^p\bigl(\norm{A-B}_{S_p}^p-T_A\bigr).
\]
Consequently, for every $\eta\ge0$,
\[
 \norm{A-B}_{S_p}^p\le(1+\eta)T_A
 \quad\Longrightarrow\quad
 \norm{f(A)-C}_{S_p}^p\le(1+2\eta)T_f.
\]
For every $\eps\ge0$, the norm-relative implication is
\[
 \norm{A-B}_{S_p}\le(1+\eps)T_A^{1/p}
 \quad\Longrightarrow\quad
 \norm{f(A)-C}_{S_p}\le(1+2\eps)T_f^{1/p}.
\]
For each fixed $p\in[1,\infty)$, both universal relative-excess constants two are optimal, even for operator-monotone power functions.
\end{theorem}

\section{Proof}
For $0\le x<y$ with $y\ge\tau$, subhomogeneity and continuity give
\[
 0\le f(y)-f(x)\le\frac{f(y)}{y}(y-x)\le c(y-x).
\]
Also $f(x)\ge cx$ for $x\le\tau$ and $f(x)\le cx$ for $x\ge\tau$. If $f(\tau)=0$, these inequalities and monotonicity force $f$ to vanish identically, and the theorem is immediate.

\begin{lemma}\label{lem:scalar}
For all $a,b\ge0$,
\[
 |f(a)-f(b)|^p\le c^p\bigl(|a-b|^p+(\tau^p-a^p)_+\bigr).
\]
\end{lemma}
\begin{proof}
If $\max(a,b)\ge\tau$, the pairwise Lipschitz inequality above suffices.
If $0\le b\le a\le\tau$ and $a>0$, subhomogeneity gives
\[
 f(a)-f(b)\le f(\tau)(1-b/a).
\]
Writing $q=1-b/a\in[0,1]$, we have
\[
 |a-b|^p+\tau^p-a^p-\tau^pq^p
 =(\tau^p-a^p)(1-q^p)\ge0,
\]
which proves this case.
If $0\le a\le b\le\tau$ and $b>0$, then
\[
 f(b)-f(a)\le f(\tau)(1-a/b).
\]
With $q=1-a/b\in[0,1]$, the required inequality follows from
\[
 (\tau^p-b^p)q^p\le\tau^p-b^p\le\tau^p-a^p,
 \qquad |a-b|^p=b^pq^p.
\]
The case $a=b=0$ is immediate; the other zero endpoints are included by continuity.
\end{proof}

\begin{proof}[Proof of Theorem~\ref{thm:commuting}]
Simultaneous diagonalization gives the exact decomposition
\[
 \delta_A:=\norm{A-B}_{S_p}^p-T_A=L_A+R_A,
\]
where
\[
 L_A=\sum_{i\le k}a_i^p-\sum_{i\in S}a_i^p,\qquad
 R_A=\sum_{i\in S}|a_i-b_i|^p.
\]
Define $L_f,R_f$ by the analogous formulas with $f(a_i),f(b_i)$; then
$\delta_f:=\norm{f(A)-C}_{S_p}^p-T_f=L_f+R_f$.

Match the selected indices outside $\{1,\ldots,k\}$ with the omitted indices in $\{1,\ldots,k\}$. Each matched pair consists of values $x\le\tau\le y$. Since $f(y)\le cy$ and $f(x)\ge cx$,
\[
 0\le f(y)^p-f(x)^p\le c^p(y^p-x^p).
\]
Summing over these pairs gives $0\le L_f\le c^pL_A$. Since every omitted $y$ is at least $\tau$, the same matching also gives
\[
 L_A\ge\sum_{i\in S}(\tau^p-a_i^p)_+.
\]
Lemma~\ref{lem:scalar} now yields
\[
 R_f\le c^p(R_A+L_A),
 \qquad
 \delta_f\le c^p(R_A+2L_A)\le2c^p\delta_A.
\]
For the relative implications, use $c^pT_A\le T_f$. Finally, convexity of $t\mapsto t^p$ gives
\[
 (1+2\eps)^p\ge2(1+\eps)^p-1,
\]
which converts the bound on the $p$th power into the norm-relative bound.
\end{proof}

\section{Sharpness and scope}
Fix $p\in[1,\infty)$. Let $N\ge1$, $0<b<1$, $0<r<1$, and take
\[
 A=I_{N+1}\oplus[0],\qquad k=1,\qquad
 B=b\,e_{N+2}e_{N+2}^T,\qquad f(x)=x^r.
\]
Here $\tau=c=1$, $T_A=T_f=N$, and
\[
 \norm{A-B}_{S_p}^p=N+1+b^p,\qquad
 \norm{f(A)-C}_{S_p}^p=N+1+b^{pr}.
\]
The ratio of $p$th-power excesses is
\[
 \frac{1+b^{pr}}{1+b^p}.
\]
Taking $b=e^{-L}$ and $r=L^{-2}$ as $L\to\infty$ makes this ratio tend to two. With the simultaneous choice $N=L$ for positive integers $L$, the ratio of norm-relative excesses is
\[
 \frac{(1+(1+b^{pr})/N)^{1/p}-1}
 {(1+(1+b^p)/N)^{1/p}-1}\longrightarrow2,
\]
using $(1+x/N)^{1/p}-1\sim x/(pN)$ for bounded positive $x$. For $p=1$, the two ratios already coincide for every $N$.

These powers are operator monotone. Indeed, for $0<r<1$,
\[
 x^r=\frac{\sin(\pi r)}{\pi}
 \int_0^\infty t^{r-1}x(x+t)^{-1}\,dt.
\]
Spectral calculus extends this identity to positive semidefinite matrices, and for $t>0$ the map
$X\mapsto X(X+tI)^{-1}=I-t(X+tI)^{-1}$ preserves their order. Integrating proves the assertion, including singular matrices.

If $\tau=0$, the relative-error premise forces $A=B$. The selected space then contains $\range A$, and $f(A)-C$ has exactly $n-k$ remaining eigenvalues equal to $f(0)$. Its Schatten $p$-error is the optimal tail. The relative conclusions therefore hold in this case too; no scaled inequality with an undefined quotient $f(0)/0$ is asserted.

For $p=2$, a stronger result allowing noncommuting matrices is proved in the companion Frobenius manuscript. For $p=1$, the present argument proves a sharp constant-two bound only with the specified simultaneous eigenbasis and selected projector. It does not settle existence of a universal constant for arbitrary noncommuting pairs. When this result is applied to a truncation $B=\widehat A_k$, the selected projector must preserve the stated common basis, including within repeated eigenspaces and at padded zero eigenvalues.

\end{document}

% END REVIEWED BODY
